%%%%%%%%%%%%%%%%%
% Ceci est un modèle de CV créé à l'aide de altacv.cls
% (v1.7.2, 28 août 2024) écrit par LianTze Lim (liantze@gmail.com). Compile avec pdfLaTeX, XeLaTeX et LuaLaTeX.
%
%% Il peut être distribué et/ou modifié sous les
%% conditions de la LaTeX Project Public License, soit la version 1.3
%% de cette licence ou (à votre choix) toute version ultérieure.
%% La dernière version de cette licence est disponible sur
%%    http://www.latex-project.org/lppl.txt
%% et la version 1.3 ou ultérieure fait partie de toutes les distributions de LaTeX
%% version 2003/12/01 ou ultérieure.
%%%%%%%%%%%%%%%%

%% Utilisez l'option "normalphoto" si vous souhaitez une photo normale au lieu d'une photo recadrée en cercle
% \documentclass[10pt,a4paper,normalphoto]{altacv}

\documentclass[8pt,a4paper,ragged2e,withhyper]{../_shared/altacv}
%% AltaCV utilise les packages fontawesome5 et simpleicons.
%% Voir http://texdoc.net/pkg/fontawesome5 et http://texdoc.net/pkg/simpleicons pour la liste complète des symboles.

% Modifiez la mise en page si nécessaire
\geometry{left=1.25cm,right=1.25cm,top=.5cm,bottom=1.cm,columnsep=1.2cm}

% Le package paracol permet de mettre en colonnes parallèles
\usepackage{paracol}
\usepackage{fontawesome5}
\usepackage{hyperref}

% Changer les couleurs si nécessaire
\definecolor{SlateGrey}{HTML}{2E2E2E}
\definecolor{LightGrey}{HTML}{666666}
\definecolor{DarkPastelRed}{HTML}{8AC0D9}%{450808}
\definecolor{PastelRed}{HTML}{00B0DA}%{8F0D0D}
\definecolor{GoldenEarth}{HTML}{B4AD0C}%{E7D192}
\colorlet{name}{black}
\colorlet{tagline}{PastelRed}
\colorlet{heading}{DarkPastelRed}
\colorlet{headingrule}{GoldenEarth}
\colorlet{subheading}{PastelRed}
\colorlet{accent}{PastelRed}
\colorlet{emphasis}{SlateGrey}
\colorlet{body}{LightGrey}

% Modifier certaines polices, si nécessaire
\renewcommand{\namefont}{\Huge\rmfamily\bfseries}
\renewcommand{\personalinfofont}{\footnotesize}
\renewcommand{\cvsectionfont}{\LARGE\rmfamily\bfseries}
\renewcommand{\cvsubsectionfont}{\large\bfseries}

% Modifier les puces pour les listes itemize et les marqueurs de notation
\renewcommand{\cvItemMarker}{{\small\textbullet}}
\renewcommand{\cvRatingMarker}{\faCircle}

% Si votre CV est dans une langue autre que l'anglais,
% il est conseillé de changer ces termes pour qu'ils apparaissent correctement en cas de copie ou d'extraction de texte depuis le PDF.
% Exemple pour le français :
% \renewcommand{\locationname}{Lieu}
% \renewcommand{\datename}{Date}

\begin{document}
\name{BOILEAU Guillaume, PhD}
\tagline{Scientifique des données}%Votre poste ou slogan ici
%% Vous pouvez ajouter plusieurs photos à gauche ou à droite
\photoL{2.8cm}{../_shared/moi}

\personalinfo{%
  \email{guillaume.boileaupro@gmail.com}
  \phone{+33 6 59 94 85 84}
  \location{Alpes-Maritimes, FRANCE}
  \linkedin{guillaume-boileau-phd-891b81265}
  \researchgate{Guillaume-Boileau-2}
  \orcid{0000-0002-3576-6968}
}

\makecvheader

%% Définir le ratio des colonnes (6:4 ici)
\columnratio{0.6}

% Début de deux colonnes parallèles
\begin{paracol}{2}

\cvsection{Expériences professionnelles}

\cvevent{\textbf{Chercheur senior}, CNES : Mission spatiale LISA (ESA/NASA)}{Laboratoire Lagrange, Observatoire de la Côte d'Azur, CNES}{2023 -- Présent}{Nice, France}
\begin{itemize}
\item \textbf{Optimisation du catalogue LISA : }Développement d'outils pour comparer et intégrer les résultats des modèles d'ondes gravitationnelles.
\item \textbf{Simulation du fond d’ondes gravitationnelles : }Modélisation du signal extragalactique pour améliorer les détections de LISA.
\end{itemize}
\divider
\cvevent{\textbf{Chercheur senior}, Télescope Einstein et LIGO/Virgo/KAGRA}{Université d’Anvers, Belgique}{2021 -- 2023}{Anvers, Belgique}
\begin{itemize}
\item \textbf{Pipeline MCMC : }Amélioration des détecteurs d’ondes gravitationnelles par réduction des sources de bruit (vibrations terrestres, champs magnétiques, etc.).
\item \textbf{Bruit newtonien : }Analyse de l’impact des vibrations terrestres sur la précision des détecteurs.
\item \textbf{LIGO/Virgo/KAGRA O4 : }Contribution à l’amélioration des détecteurs pour la prochaine phase d’observation.
\end{itemize}
\divider
\cvevent{\textbf{Chercheur junior}, Observations cosmologiques avec LISA : limitations dues au bruit et aux signaux astrophysiques}{ARTEMIS, Observatoire de la Côte d'Azur}{2018 -- 2021}{Nice, France}
\begin{itemize}
\item \textbf{Séparation spectrale des ondes gravitationnelles : }Analyse des signaux stochastiques pour la mission LISA.
\item \textbf{Calcul du fond galactique : }Évaluation des contributions des systèmes binaires.
\item \textbf{Recherche de bruit corrélé : }Étude des interférences potentielles dans les missions LISA et PathFinder.
\end{itemize}

\cvsection{Projets}

\cvevent{Membre du consortium LISA}{Mission spatiale ESA}{2018 -- Présent}{}
Responsable stratégique pour l’analyse des données en France (CNES). \\
\textbf{Budget : }Environ 2 milliards d’euros sur 10 ans.
\divider
\cvevent{Membre des collaborations Einstein Telescope et LIGO-Virgo-KAGRA}{Collaborations internationales}{2021 -- Présent}{}
Participation aux efforts de développement et d’amélioration des performances analytiques. \\
\textbf{Prix Nobel de Physique 2017 : }Attribué à la collaboration LIGO pour la découverte des ondes gravitationnelles.

\cvsection{Langues}
\cvskill{Français}{5}
\divider
\cvskill{Anglais}{3.5}
\divider
\cvskill{Espagnol}{2}

\end{paracol}
\end{document}
